\documentclass{article}
\usepackage[T1]{fontenc}
\usepackage{lmodern}
\usepackage{graphicx}
\usepackage{amsmath,amssymb}
\usepackage[a4paper,margin=2cm]{geometry}
\usepackage[absolute,overlay]{textpos}
\setlength{\TPHorizModule}{1mm}
\setlength{\TPVertModule}{1mm}
\usepackage[indonesian]{babel}
\usepackage{hyperref}
\setlength{\emergencystretch}{2em}
% unit: Halaman 19

\begin{document}

% === Halaman 19 (python/native) ===

Enkripsi

• Tinjau RC5 dengan ukuran blok 64 bit dan jumlah putaran r. 

• Enkripsi menggunakan kunci putaran S0, S1, …, S2r + 2 yang masing-masing 

panjangnya 32-bit. 

• Dua kunci putaran digunakan untuk setiap putaran i = 1, 2, .. r dan dua buah kunci 

putaran tambahan sebelum putaran pertama jadi seluruhnya ada 2r + 2 buah 

kunci internal). 

• Untuk melakukan enkripsi, mula-mula blok plainteks dibagi menjadi 2 bagian, A 

dan B, yang masing-masing panjangnya 32 bit. Kemudian masing-masing bagian 

dijumlahkan (dalam modulo 232) dengan S0 dan S1:



A  A + S[0]



B  B + S[1]

19

\end{document}
